% Section declaración - ECONOMÍA / ANÁLISIS ECONÓMICO
\par{

	Economista con sólida formación en análisis macroeconómico, microeconomía aplicada y econometría. Especializado en elaboración de estudios económicos, proyecciones de indicadores, análisis de coyuntura y evaluación de políticas públicas con enfoque cuantitativo.

	Experiencia en procesamiento de bases de datos económicas (INEI, BCRP, MEF), elaboración de boletines de análisis sectorial, seguimiento de variables macroeconómicas y redacción de informes técnicos. Dominio de software econométrico (Stata, EViews, R) y capacidad de comunicar hallazgos económicos complejos de manera clara y precisa.

	Busco oportunidades en áreas de estudios económicos, análisis sectorial o investigación aplicada donde pueda contribuir con mi capacidad analítica, rigor metodológico y conocimiento actualizado de la realidad económica nacional e internacional para la formulación de políticas, estrategias y recomendaciones basadas en evidencia.

	%%%%%%%%%%%%%%%%%%%%%%%%%%%%%%%%%%%%%%%%
	% VERSIÓN ALTERNATIVA - Investigación económica
	%%%%%%%%%%%%%%%%%%%%%%%%%%%%%%%%%%%%%%%%
	%Economista con especialización en investigación económica aplicada y evaluación de impacto. Experto en diseño y ejecución de estudios cuantitativos para análisis de mercados, evaluación de programas sociales y proyecciones económicas sectoriales.

	%Dominio de metodologías econométricas avanzadas: series de tiempo, datos panel, modelos de elección discreta, diferencias en diferencias y métodos cuasiexperimentales. Capacidad de trabajar con microdatos de encuestas nacionales (ENAHO, ENE, ENAPRES) y bases administrativas para investigación empírica rigurosa.

	%Busco contribuir en centros de investigación económica, think tanks o áreas de estudios donde pueda desarrollar investigaciones de frontera, publicar hallazgos en revistas especializadas y generar evidencia robusta para el diseño de políticas públicas efectivas.

	%%%%%%%%%%%%%%%%%%%%%%%%%%%%%%%%%%%%%%%%
	% VERSIÓN ALTERNATIVA - Análisis de coyuntura
	%%%%%%%%%%%%%%%%%%%%%%%%%%%%%%%%%%%%%%%%
	%Economista especializado en análisis de coyuntura económica y seguimiento de indicadores macroeconómicos. Experto en monitoreo diario de variables clave: inflación, tipo de cambio, empleo, PBI sectorial, balanza comercial y variables fiscales.

	%Capacidad demostrada en elaboración de reportes de coyuntura, notas técnicas rápidas, alertas económicas y síntesis ejecutivas para diferentes niveles de audiencia. Familiarizado con fuentes oficiales (BCRP, INEI, MEF, FED, BCE) y análisis de riesgo país.

	%Busco oportunidades en áreas de análisis macroeconómico, risk management o research donde pueda aplicar mi capacidad de síntesis, velocidad de respuesta ante cambios del entorno económico y habilidad para comunicar implicancias de eventos económicos nacionales e internacionales.

	%%%%%%%%%%%%%%%%%%%%%%%%%%%%%%%%%%%%%%%%
	% VERSIÓN ALTERNATIVA - Planificación económica
	%%%%%%%%%%%%%%%%%%%%%%%%%%%%%%%%%%%%%%%%
	%Economista con orientación en planificación económica y desarrollo regional. Experiencia en elaboración de planes de desarrollo económico local, análisis de brechas socioeconómicas, identificación de potencialidades territoriales y diseño de estrategias de desarrollo productivo.

	%Dominio de herramientas de planificación estratégica, análisis FODA, matrices de priorización y diseño de indicadores de seguimiento. Capacidad de articular actores públicos, privados y sociedad civil para construcción de visiones compartidas de desarrollo territorial.

	%Busco contribuir en gobiernos regionales, locales o agencias de desarrollo donde pueda aplicar mi formación económica y mi visión territorial para diseñar e implementar políticas que impulsen el desarrollo económico sostenible, la reducción de brechas y el bienestar de la población.

	%%%%%%%%%%%%%%%%%%%%%%%%%%%%%%%%%%%%%%%%
	% VERSIÓN ALTERNATIVA - Economía sectorial
	%%%%%%%%%%%%%%%%%%%%%%%%%%%%%%%%%%%%%%%%
	%Economista con especialización en análisis sectorial y estudios de mercado. Experto en caracterización de cadenas productivas, análisis de competitividad sectorial, estudios de oferta y demanda, y elaboración de estudios de factibilidad económica para proyectos de inversión.

	%Experiencia en análisis agrario, comercio exterior, turismo y servicios. Capacidad de realizar estudios de mercado, benchmarking sectorial, análisis de tendencias y proyecciones de demanda utilizando metodologías cuantitativas rigurosas.

	%Busco oportunidades en áreas de inteligencia comercial, desarrollo de negocios o consultoría económica donde pueda aplicar mi conocimiento sectorial profundo y mis habilidades analíticas para identificar oportunidades, evaluar viabilidad de proyectos y generar recomendaciones estratégicas basadas en evidencia.

}
